
\chapter{Matrici densità di un sistema composto}

Nella sezione precedente abbiamo costruito lo spazio di Hilbert
$\mathcal{H} = \mathcal{H}_1 \otimes \mathcal{H}_2$ del sistema composto.
Gli stati del sistema globale sono dunque gli elementi di
$\mathfrak{S}(\mathcal{H})$, e il nostro problema è capire che relazione
intercorra fra uno stato globale $\hat{\rho}$ e gli stati dei due
sottosistemi presi separatamente.

Fissate due basi ortonormali $\{\ket{a_m}\}$ di $\mathcal{H}_1$ e
$\{\ket{b_n}\}$ di $\mathcal{H}_2$, un generico $\hat{\rho} \in
\mathfrak{S}(\mathcal{H})$ è individuato dai suoi elementi di matrice
$\braket{a_m, b_n | \hat{\rho} | a_{m'}, b_{n'}}$ sulla base prodotto.

\subsection{Traccia parziale e stati ridotti}

Consideriamo un'osservabile che appartenga \emph{esclusivamente} al primo
sottosistema. Come visto nella sezione precedente, il suo operatore ha la
forma $\hat{R} \otimes \hat{\mathbb{1}}$. Il suo valor medio è
%
\begin{align}
    \bigl\langle \hat{R} \otimes \hat{\mathbb{1}} \bigr\rangle
    &= \operatorname{tr} \bigl[ \hat{\rho} \, (\hat{R} \otimes
       \hat{\mathbb{1}}) \bigr] \nonumber \\
    &= \sum_{m,n} \sum_{m',n'}
       \braket{a_m, b_n | \hat{\rho} | a_{m'}, b_{n'}}
       \braket{a_{m'}, b_{n'} | \hat{R} \otimes \hat{\mathbb{1}} | a_m, b_n}
       \nonumber \\
    &= \sum_{m,m'} \biggl( \sum_n
       \braket{a_m, b_n | \hat{\rho} | a_{m'}, b_n} \biggr)
       \braket{a_{m'} | \hat{R} | a_m}
    \label{eq:media-locale}
\end{align}
%
dove nell'ultimo passaggio si è usato
$\braket{a_{m'}, b_{n'} | \hat{R} \otimes \hat{\mathbb{1}} | a_m, b_n}
= \braket{a_{m'} | \hat{R} | a_m} \, \delta_{n'n}$.

Il punto notevole è che la quantità fra parentesi dipende solo da
$\hat{\rho}$ e non da $\hat{R}$: definisce un operatore su
$\mathcal{H}_1$ che contiene tutta l'informazione necessaria a calcolare
le medie locali. Lo chiamiamo \emph{matrice densità ridotta}, e
l'operazione che lo produce \emph{traccia parziale}%
\mkeyword{traccia parziale}.
%
\begin{tcolorbox}[colback=black!5!white, colframe=black!45!white, boxrule=0.8pt]
\begin{equation}
    \hat{\rho}^{(1)} \equiv \operatorname{tr}_2 \hat{\rho} \, , \qquad
    \braket{a_m | \hat{\rho}^{(1)} | a_{m'}} = \sum_n
    \braket{a_m, b_n | \hat{\rho} | a_{m'}, b_n}
    \label{eq:traccia-parziale}
\end{equation}
\end{tcolorbox}
%
e analogamente $\hat{\rho}^{(2)} = \operatorname{tr}_1 \hat{\rho}$,
sommando sull'indice del primo sottosistema. Con questa definizione la
\eqref{eq:media-locale} si scrive
%
\begin{equation}
    \bigl\langle \hat{R} \otimes \hat{\mathbb{1}} \bigr\rangle
    = \operatorname{tr}_1 \bigl( \hat{\rho}^{(1)} \hat{R} \bigr)
    \label{eq:media-ridotta}
\end{equation}
%
che è formalmente identica alla regola per un sistema singolo. Va
sottolineato che $\hat{\rho}^{(1)}$ non è stato postulato ma
\emph{dedotto}: è l'unico operatore su $\mathcal{H}_1$ che riproduca le
medie di tutte le osservabili locali.

Perché la definizione sia sensata occorre però verificare che
$\hat{\rho}^{(1)}$ sia effettivamente uno stato.
%
\begin{tcolorbox}[colback=black!5!white, colframe=black!45!white, boxrule=0.8pt]
\textbf{Proposizione.} Se $\hat{\rho} \in \mathfrak{S}(\mathcal{H})$,
allora $\hat{\rho}^{(1)} \in \mathfrak{S}(\mathcal{H}_1)$.
\end{tcolorbox}
%
\begin{proof}
Verifichiamo le due condizioni che definiscono $\mathfrak{S}$.

Per la traccia, sommando gli elementi diagonali della
\eqref{eq:traccia-parziale},
%
\begin{equation*}
    \operatorname{tr}_1 \hat{\rho}^{(1)} = \sum_m
    \braket{a_m | \hat{\rho}^{(1)} | a_m} = \sum_{m,n}
    \braket{a_m, b_n | \hat{\rho} | a_m, b_n}
    = \operatorname{tr} \hat{\rho} = 1
\end{equation*}
%
dove si è riconosciuta la traccia completa sulla base prodotto.

Per la positività, sia $\ket{u} = \sum_m u_m \ket{a_m}$ un vettore
arbitrario di $\mathcal{H}_1$. Allora
%
\begin{align*}
    \braket{u | \hat{\rho}^{(1)} | u}
    &= \sum_{m,m'} \overline{u_m} \, u_{m'} \sum_n
       \braket{a_m, b_n | \hat{\rho} | a_{m'}, b_n} \\
    &= \sum_n \Bigl( \sum_m \overline{u_m} \bra{a_m, b_n} \Bigr)
       \hat{\rho} \Bigl( \sum_{m'} u_{m'} \ket{a_{m'}, b_n} \Bigr) \\
    &= \sum_n \braket{u \otimes b_n | \hat{\rho} | u \otimes b_n}
    \ \geq \ 0
\end{align*}
%
poiché ciascun addendo è non negativo per la positività di $\hat{\rho}$.
L'hermitianità segue dalla positività su uno spazio di Hilbert
complesso.
\end{proof}

\subsection{Decomposizione di uno stato generico}

Disponendo delle due ridotte, è naturale chiedersi se esse determinino
$\hat{\rho}$. La risposta è negativa, e conviene renderla quantitativa.

Introduciamo sullo spazio degli operatori il prodotto scalare di
Hilbert--Schmidt%
\mkeyword{prodotto di Hilbert--Schmidt}
%
\begin{equation}
    \bigl( \hat{A}, \hat{B} \bigr)_{\mathrm{HS}} \equiv
    \operatorname{tr} \bigl( \hat{A}^\dagger \hat{B} \bigr)
    \label{eq:hilbert-schmidt}
\end{equation}
%
e scegliamo in $\mathcal{L}_{\mathrm{h}}(\mathcal{H}_1)$ una base
ortonormale $\{\hat{\tau}^{(1)}_\mu\}_{\mu = 0}^{d_1^2 - 1}$ rispetto a
tale prodotto, con la convenzione
%
\begin{equation}
    \hat{\tau}^{(1)}_0 = \frac{\hat{\mathbb{1}}_1}{\sqrt{d_1}} \, ,
    \qquad
    \operatorname{tr} \hat{\tau}^{(1)}_\mu = 0 \quad \text{per } \mu \geq 1
    \label{eq:base-operatori}
\end{equation}
%
sempre possibile, dato che gli operatori hermitiani a traccia nulla sono
il complemento ortogonale dell'identità. Procediamo allo stesso modo per
il secondo sottosistema.

Poiché $\mathcal{L}(\mathcal{H}_1 \otimes \mathcal{H}_2) =
\mathcal{L}(\mathcal{H}_1) \otimes \mathcal{L}(\mathcal{H}_2)$, i prodotti
$\hat{\tau}^{(1)}_\mu \otimes \hat{\tau}^{(2)}_\nu$ formano una base
ortonormale per gli operatori sul sistema composto, e ogni stato si
espande in modo unico come
%
\begin{equation}
    \hat{\rho} = \sum_{\mu, \nu} c_{\mu\nu} \,
    \hat{\tau}^{(1)}_\mu \otimes \hat{\tau}^{(2)}_\nu \, ,
    \qquad
    c_{\mu\nu} = \operatorname{tr} \bigl[ \hat{\rho} \,
    ( \hat{\tau}^{(1)}_\mu \otimes \hat{\tau}^{(2)}_\nu ) \bigr]
    \in \mathbb{R}
    \label{eq:espansione-rho}
\end{equation}

Calcoliamo ora le ridotte in questa rappresentazione. Usando
$\operatorname{tr}_2 (\hat{A} \otimes \hat{B}) = \hat{A}
\operatorname{tr} \hat{B}$ insieme alla \eqref{eq:base-operatori},
sopravvive il solo termine $\nu = 0$:
%
\begin{equation}
    \hat{\rho}^{(1)} = \sqrt{d_2} \sum_\mu c_{\mu 0} \,
    \hat{\tau}^{(1)}_\mu \, ,
    \qquad
    \hat{\rho}^{(2)} = \sqrt{d_1} \sum_\nu c_{0 \nu} \,
    \hat{\tau}^{(2)}_\nu
    \label{eq:ridotte-coefficienti}
\end{equation}

Disponendo i coefficienti in una matrice $c$, le ridotte leggono dunque
soltanto la \emph{prima riga} e la \emph{prima colonna}. Tutti i
coefficienti con $\mu \geq 1$ \emph{e} $\nu \geq 1$ sono invisibili a
entrambe: sono esattamente l'informazione che lo stato globale possiede e
che la coppia degli stati parziali non registra.

Il conteggio è immediato. Uno stato su $\mathcal{H}$ è individuato da
$(d_1 d_2)^2 - 1$ parametri reali, le due ridotte ne portano
rispettivamente $d_1^2 - 1$ e $d_2^2 - 1$, e restano
%
\begin{equation}
    (d_1 d_2)^2 - 1 - (d_1^2 - 1) - (d_2^2 - 1)
    = (d_1^2 - 1)(d_2^2 - 1)
    \label{eq:conteggio-correlazioni}
\end{equation}
%
parametri di pura correlazione. Per due qubit si tratta di $9$ parametri
su $15$: più della metà dell'informazione contenuta nello stato globale
non è ricostruibile dai sottosistemi.

Possiamo allora isolare la parte di $\hat{\rho}$ che le ridotte non
vedono definendo l'\emph{operatore di correlazione}%
\mkeyword{operatore di correlazione}
%
\begin{tcolorbox}[colback=black!5!white, colframe=black!45!white, boxrule=0.8pt]
\begin{equation}
    \hat{\rho} = \hat{\rho}^{(1)} \otimes \hat{\rho}^{(2)} + \hat{\chi}
    \, , \qquad
    \operatorname{tr}_1 \hat{\chi} = \operatorname{tr}_2 \hat{\chi} = 0
    \label{eq:decomposizione-correlazione}
\end{equation}
\end{tcolorbox}
%
Le due proprietà di traccia si verificano subito: ad esempio
%
\begin{equation*}
    \operatorname{tr}_2 \hat{\chi} = \operatorname{tr}_2 \hat{\rho}
    - \operatorname{tr}_2 \bigl( \hat{\rho}^{(1)} \otimes
    \hat{\rho}^{(2)} \bigr)
    = \hat{\rho}^{(1)} - \hat{\rho}^{(1)}
    \operatorname{tr} \hat{\rho}^{(2)} = 0
\end{equation*}
%
L'operatore $\hat{\chi}$ è hermitiano a traccia nulla, ma in generale
\emph{non} è positivo, e dunque non è uno stato: è una correzione, non
una matrice densità%
\footnote{Non potrebbe esserlo: un operatore positivo a traccia nulla è
l'operatore nullo.}.

Il significato di $\hat{\chi}$ emerge calcolando una correlazione fra
osservabili locali. Detti $\hat{R}_1 = \hat{R}^{(1)} \otimes
\hat{\mathbb{1}}$ e $\hat{R}_2 = \hat{\mathbb{1}} \otimes \hat{R}^{(2)}$,
%
\begin{equation}
    \bigl\langle \hat{R}_1 \hat{R}_2 \bigr\rangle
    - \bigl\langle \hat{R}_1 \bigr\rangle
      \bigl\langle \hat{R}_2 \bigr\rangle
    = \operatorname{tr} \bigl[ \hat{\chi} \,
      ( \hat{R}^{(1)} \otimes \hat{R}^{(2)} ) \bigr]
    \label{eq:correlazione-chi}
\end{equation}
%
dunque $\hat{\chi}$ è precisamente l'operatore che trasporta tutte le
correlazioni fra i due sottosistemi.

\paragraph{Esempio: due qubit.}
Scegliendo come base degli operatori le matrici di Pauli normalizzate,
$\hat{\tau}_0 = \hat{\mathbb{1}} / \sqrt{2}$ e $\hat{\tau}_j =
\hat{\sigma}_j / \sqrt{2}$, la \eqref{eq:espansione-rho} assume la forma
%
\begin{equation}
    \hat{\rho} = \frac{1}{4} \Bigl(
    \hat{\mathbb{1}} \otimes \hat{\mathbb{1}}
    + \sum_j a_j \, \hat{\sigma}_j \otimes \hat{\mathbb{1}}
    + \sum_k b_k \, \hat{\mathbb{1}} \otimes \hat{\sigma}_k
    + \sum_{j,k} T_{jk} \, \hat{\sigma}_j \otimes \hat{\sigma}_k \Bigr)
    \label{eq:forma-fano}
\end{equation}
%
dove $\mathbf{a}$ e $\mathbf{b}$ sono i vettori di polarizzazione dei due
qubit, che determinano le ridotte, e $T$ è la \emph{matrice di
correlazione}. I conti tornano: $3 + 3 + 9 = 15$.

Il caso del singoletto è istruttivo. Entrambe le ridotte sono
$\hat{\rho}^{(1)} = \hat{\rho}^{(2)} = \hat{\mathbb{1}}/2$, cioè
$\mathbf{a} = \mathbf{b} = \mathbf{0}$, mentre $T_{jk} = -\delta_{jk}$:
%
\begin{equation}
    \ket{\Psi^-}\bra{\Psi^-} = \frac{1}{4} \Bigl(
    \hat{\mathbb{1}} \otimes \hat{\mathbb{1}}
    - \sum_j \hat{\sigma}_j \otimes \hat{\sigma}_j \Bigr)
    \label{eq:singoletto-fano}
\end{equation}
%
Tutta l'informazione sullo stato risiede nella matrice di correlazione e
nessuna nelle ridotte: è la realizzazione estrema dell'affermazione da
cui siamo partiti.

\subsection{Stati non correlati}

Il caso in cui $\hat{\chi}$ si annulla merita un nome%
\mkeyword{stato non correlato}.
%
\begin{tcolorbox}[colback=black!5!white, colframe=black!45!white, boxrule=0.8pt]
Uno stato $\hat{\rho} \in \mathfrak{S}(\mathcal{H}_1 \otimes
\mathcal{H}_2)$ si dice \emph{non correlato} se
%
\begin{equation}
    \hat{\rho} = \hat{\rho}^{(1)} \otimes \hat{\rho}^{(2)}
    \label{eq:non-correlato}
\end{equation}
\end{tcolorbox}
%
Nella rappresentazione \eqref{eq:espansione-rho} questa condizione si
legge in modo trasparente: $\hat{\rho}$ è non correlato se e solo se
$c_{\mu\nu} = r_\mu s_\nu$, cioè se e solo se la matrice dei
coefficienti $c$ ha \emph{rango} $1$. È lo stesso criterio che
caratterizza i vettori fattorizzabili, trasportato dal livello dei
vettori a quello degli operatori.

La caratterizzazione operativamente utile è però un'altra.
%
\begin{tcolorbox}[colback=black!5!white, colframe=black!45!white, boxrule=0.8pt]
\textbf{Teorema.} $\hat{\rho}$ è non correlato se e solo se per ogni
coppia di osservabili locali $\hat{R}_1 = \hat{R}^{(1)} \otimes
\hat{\mathbb{1}}$ e $\hat{R}_2 = \hat{\mathbb{1}} \otimes \hat{R}^{(2)}$
vale
%
\begin{equation}
    \bigl\langle \hat{R}_1 \hat{R}_2 \bigr\rangle
    = \bigl\langle \hat{R}_1 \bigr\rangle
      \bigl\langle \hat{R}_2 \bigr\rangle
    \label{eq:fattorizzazione-medie}
\end{equation}
\end{tcolorbox}
%
\begin{proof}
Si osservi anzitutto che $[\hat{R}_1, \hat{R}_2] = 0$, cosicché il primo
membro è ben definito e reale.

$(\Rightarrow)$ Il prodotto dei due operatori è $\hat{R}^{(1)} \otimes
\hat{R}^{(2)}$, e usando
$(\hat{A} \otimes \hat{B})(\hat{C} \otimes \hat{D}) = \hat{A}\hat{C}
\otimes \hat{B}\hat{D}$ insieme a
$\operatorname{tr} (\hat{A} \otimes \hat{B}) =
\operatorname{tr}_1 \hat{A} \cdot \operatorname{tr}_2 \hat{B}$,
%
\begin{equation*}
    \bigl\langle \hat{R}_1 \hat{R}_2 \bigr\rangle
    = \operatorname{tr} \bigl[ ( \hat{\rho}^{(1)} \hat{R}^{(1)} )
      \otimes ( \hat{\rho}^{(2)} \hat{R}^{(2)} ) \bigr]
    = \operatorname{tr}_1 \bigl( \hat{\rho}^{(1)} \hat{R}^{(1)} ) \cdot
      \operatorname{tr}_2 \bigl( \hat{\rho}^{(2)} \hat{R}^{(2)} )
\end{equation*}
%
che per la \eqref{eq:media-ridotta} è il secondo membro della
\eqref{eq:fattorizzazione-medie}.

$(\Leftarrow)$ L'ipotesi afferma che
%
\begin{equation*}
    \operatorname{tr} \bigl[ \hat{\rho} \, ( \hat{R}^{(1)} \otimes
    \hat{R}^{(2)} ) \bigr]
    = \operatorname{tr} \bigl[ ( \hat{\rho}^{(1)} \otimes
    \hat{\rho}^{(2)} ) ( \hat{R}^{(1)} \otimes \hat{R}^{(2)} ) \bigr]
\end{equation*}
%
per ogni coppia $\hat{R}^{(1)}, \hat{R}^{(2)}$. Ma gli operatori prodotto
generano $\mathcal{L}(\mathcal{H}_1 \otimes \mathcal{H}_2)$, e il
prodotto di Hilbert--Schmidt \eqref{eq:hilbert-schmidt} è non degenere:
due operatori con la stessa traccia contro ogni elemento di un insieme
generatore coincidono. Dunque $\hat{\rho} = \hat{\rho}^{(1)} \otimes
\hat{\rho}^{(2)}$.
\end{proof}

Uno stato non correlato è dunque quello per cui le osservabili dei due
sottosistemi sono \emph{statisticamente indipendenti}, nel senso
ordinario della teoria della probabilità. Equivalentemente, è il solo
caso in cui la coppia $(\hat{\rho}^{(1)}, \hat{\rho}^{(2)})$ determina
$\hat{\rho}$, e dunque il solo caso in cui lo stato del sistema globale
non è nulla di più degli stati dei suoi sottosistemi.

Due avvertenze per chiudere.

La prima è che la contronominale del teorema è la direzione utile nelle
applicazioni: l'osservazione di una qualunque correlazione non nulla
certifica che lo stato globale non è fattorizzato.

La seconda riguarda il lessico, ed è una distinzione che dovremo
riprendere. «Non correlato» \emph{non} è sinonimo di «non entangled»:
esistono stati correlati che restano descrivibili come miscele statistiche
di stati prodotto, e in cui la correlazione è di natura puramente
classica. Un esempio è
%
\begin{equation}
    \hat{\rho} = \tfrac{1}{2} \bigl(
    \ket{{\uparrow}{\uparrow}} \bra{{\uparrow}{\uparrow}}
    + \ket{{\downarrow}{\downarrow}} \bra{{\downarrow}{\downarrow}}
    \bigr)
    \label{eq:correlazione-classica}
\end{equation}
%
che ha $\hat{\chi} \neq 0$ eppure non contiene alcun entanglement: è la
situazione di due monete truccate che danno sempre lo stesso esito. La
nozione di separabilità, strettamente più debole di quella di assenza di
correlazione, sarà oggetto della prossima sezione.

% \newpage

% Definendo l'insieme dei possibili stati di un sistema come\mkeyword{insieme degli stati}
% %
% \begin{equation}
%     \mathfrak{S}\mathcal{H} = \{ \hat{\rho} \in \mathcal{L}(\mathcal{H}) \ ; \ \text{tr}(\hat{\rho}) = 1 \ , \ \hat{\rho} \geq 0\ \ , \ \hat{\rho} = \hat{\rho}^{\dagger}\} \subset \mathcal{L}(\mathcal{H})
% \end{equation}
% dove $\mathcal{L}(\mathcal{H})$ è l'insieme degli operatori lineari sullo spazio di Hilbert $\mathcal{H}$. 